\documentclass[10pt,a4paper]{amsart}
\usepackage[margin=19mm]{geometry}
\usepackage{amsmath,amssymb,mathtools}
\usepackage[scr=rsfs,cal=euler]{mathalfa}
\usepackage{xfrac}
\usepackage[unicode,hidelinks,bookmarks=false]{hyperref}
\newcommand{\ie}{i.e.\ }
\newcommand{\cf}{cf.\ }
\newcommand{\resp}{respectively\ }
\newcommand{\Subsection}{\subsection}
\providecommand{\nobreakdash}{\nobreak\hbox{-}}
\setlength{\emergencystretch}{2em}
\multlinegap0pt
\usepackage[shortlabels]{enumitem}
\usepackage{tikz}
\usepackage{xcolor}
\usepackage{xfrac}
\usepackage{amssymb}
\usepackage[normalem]{ulem}
\newtheorem{definition}{Definition}
\newtheorem{theorem}{Theorem}
\newtheorem{proposition}{Proposition}
\newcommand{\R}{\mathbb R}
\newcommand{\T}{\mathbb T}
\newcommand{\Z}{\mathbb Z}
\title{Nontrivial Integrable Weak Stationary Solutions to Active Scalar Equations with Non-Odd Drift}
\author{Nicholas Gismondi}
\date{Source: arXiv:2601.14592v1 (CC BY 4.0)}
\begin{document}
\maketitle

\noindent\emph{Source and licence.} Nontrivial Integrable Weak Stationary Solutions to Active Scalar Equations with Non-Odd Drift. Version arXiv:2601.14592v1 (CC BY 4.0), distributed by arXiv under the Creative Commons Attribution 4.0 International licence, http://creativecommons.org/licenses/by/4.0/, as recorded in the arXiv licence metadata for this exact version and shown on the abstract page https://arxiv.org/abs/2601.14592v1. Full licence notice: \texttt{licenses/arxiv-2601.14592-CC-BY-4.0.txt}.\par\smallskip

\noindent\textit{Editorial scope.} Counted unit: the article's proposition prop:ind (source0.tex lines 656--684). The article's complete original proof of it is delivered with it (source0.tex lines 686--729, one \texttt{proof} environment of 44 lines). No mathematics is written, repaired or re-derived: the statement and the proof are the article's own lines, in the article's own notation, and no retained line is substituted or re-flowed.  Retained uncounted as the source-local context the proof needs: lines 256--261; lines 331--341; lines 345--351; lines 355--366; lines 424--438; lines 610--615; lines 744--744.  Nothing was omitted from any retained span, so the package carries no omission record.  No retained line contains a citation, so the wrapper carries no bibliography and no bibliography entry.  No retained line contains a figure environment, an image-inclusion command or drawing code, so no image is delivered and the package declares no asset dependency.  Editorial formatting only, in addition: an AMS \texttt{amsart} preamble that restates the article's own declarations and macros used by the retained lines, namely the environments \texttt{definition}, \texttt{theorem}, \texttt{proposition} on separate counters, together with the standard packages those lines need.  The retained environments print in this document as Definition 1, Theorem 1, Proposition 1 in order of appearance, whereas the article prints the same items with the numbering of its own full document.  The complete native source of the article as distributed in the arXiv e-print for this exact version is bundled unchanged as \texttt{sources/193100/source0.tex}.
\begin{equation}\label{eq:stat_eqn}
\begin{cases}
    \nabla \cdot (\theta u) + \Lambda^\gamma \theta = 0\\
    u = T\theta
\end{cases}
\end{equation}

\begin{definition}[\textbf{Paraproducts in $\dot H^s(\T^d)$}]\label{def:paras}
    Let $f,g$ be distributions, so that $\mathbb{P}_{2^j}(f), \mathbb{P}_{2^{j'}}(g)$ are well-defined for $j,j'\geq 0$ (see Definition \ref{def:projs}).  We say that $\mathbb{P}_{\not = 0}(fg)$ is well-defined as a paraproduct in $\dot{H}^s(\T^d)$ for some $s \in \R$ if
    $$
        \sum_{j,j' \geq 0} \left\Vert \mathbb{P}_{\not = 0}\left(\mathbb{P}_{2^j}(f) \mathbb{P}_{2^{j'}}(g)\right) \right\Vert_{\dot{H}^s} < \infty \, .
    $$
    Then we define
    $$
        \mathbb{P}_{\not = 0}(fg) = \sum_{j,j' \geq 0} \mathbb{P}_{\not = 0} \left(\mathbb{P}_{2^j}(f) \mathbb{P}_{2^{j'}}(g)\right) \, ,
    $$
    since the right-hand side is an absolutely summable series in $\dot H^s(\T^d)$.
\end{definition}

\begin{definition}[\textbf{Weak paraproduct solutions to ~\eqref{eq:stat_eqn}}] \label{def:weak_para_soln}
    If $\theta \in \dot{H}^s$, $s < 0$, and $u = T \theta \in \dot{H}^s$, we say $u$ and $\theta$ form a weak paraproduct solution to the stationary active scalar equation if there is $s' \in \R$ such that $\mathbb{P}_{\not = 0}(\theta u)$ is well defined as a paraproduct in $\dot{H}^{s'}$ in the sense of the previous definition and
    $$
        \left\langle \theta, \Lambda^{\gamma} \phi \right\rangle_{\dot{H}^{s},\dot{H}^{-s}} - \left\langle \mathbb{P}_{\not= 0}(\theta u), \nabla \phi \right\rangle_{\dot{H}^{s'}, \dot{H}^{-s'}} = 0
    $$
    for all smooth $\phi$.
\end{definition}

\begin{theorem}\label{thm:main}
    Given any $0 < \gamma \leq 2$ there exist
    $$
    u,\theta \in \bigcap_{0 < \epsilon < 1} \left(\dot{B}_{\infty,\infty}^{-\epsilon}(\T^d) \cap L^{2-\epsilon}(\T^d)\right) \setminus \{0\}
    $$
    such that
    \begin{enumerate}
        \item [(a)] $u = T\theta$;
        \item [(b)] $\mathbb{P}_{\not=0}(\theta u)$ exists as a paraproduct in $\dot{H}^{-s}(\T^d)$ for all $s > d/2 + \max(\gamma - 1, 0)$;
        \item [(c)] $u$ and $\theta$ solve Equation ~\eqref{eq:stat_eqn} in the sense of Definition \ref{def:weak_para_soln}.
    \end{enumerate}
\end{theorem}

\begin{definition}[\textbf{Littlewood-Paley projectors}]\label{def:projs}
    There exists $\varphi \colon \R^{d}\to[0,1]$, smooth, radially symmetric, and compactly supported in $\{\xi : 6/7 \leq |\xi|\leq 2\}$ such that $\varphi(\xi) = 1$ on $\{\xi : 1 \leq |\xi| \leq 12/7\}$,
    \begin{equation}
        \sum_{j\geq 0}\varphi(2^{-j}\xi)=1 \hspace{0.25cm} \text{ for all } \hspace{0.25cm} |\xi|\geq 1, \notag
    \end{equation}
    and $\operatorname{supp}\varphi_{j}\cap \operatorname{supp} \varphi_{j'}=\emptyset$ for all $|j-j'|\geq 2$, where $\varphi_j(\cdot) = \varphi(2^{-j}\cdot)$. We define the projection of a function $f$ on its $0$-mode by
    \begin{equation}
        \mathbb{P}_{=0}(f)=\int_{\T^d} f(x)\, dx, \notag
    \end{equation}
    and the projection on the $j^{\rm th}$ shell by
    \begin{equation}
        \mathbb{P}_{2^j}(f)(x)=\sum_{k\in \Z^d}\hat{f}(k)\varphi_{j}(k)e^{2\pi ik\cdot x} \, . \notag
    \end{equation}
    We also define $\mathbb{P}_{\neq 0}f:=(\operatorname{Id}-\mathbb{P}_{=0})f$, and we also denote by $\hat{K}_{\simeq 1}$ a smooth radially symmetric bump function with support in the ball $\{\xi: |\xi| < r\} $, which also satisfies $\hat{K}_{\simeq 1}(\xi)=1$ on the smaller ball $\{\xi: |\xi| \leq \frac{1}{2}r\}$ for $r \ll 1$ is a power of $2$ whose exact size will be specified in Section \ref{sec:param}. Then let $\mathbb{P}_{\leq \lambda}$ be the convolution operator that has $\hat{K}_{\simeq 1}\left(\frac{\xi}{\lambda}\right)$ as its Fourier multiplier.
\end{definition}

\begin{equation}\label{eq:relaxed}
\begin{split}
    \nabla \cdot (\theta u) + \Lambda^\gamma \theta &= \nabla \cdot R\\
    u &= T\theta
    \end{split}
\end{equation}

\subsection{Parameters}\label{sec:param}

\begin{proposition}\label{prop:ind}
    We make the following inductive assumptions about $(\theta_q,u_q,R_q)$:
    \begin{enumerate}
        \item\label{i:1} $\theta_q$ has zero mean and $u_q$ is divergence free;
        \item\label{i:2} $(\theta_q,u_q,R_q)$ is a smooth solution to the relaxed active scalar equation given by Equation ~\eqref{eq:relaxed};
        \item\label{i:3} $\Vert R_q \Vert_{\dot{H}^{-s}} < 2^{-q}$;
        \item\label{i:4} There is a constant $C(\alpha,p) > 0$ depending only on $\alpha < 0$ and $1 \leq p < 2$ such that for all $q' \leq q$ we have
        $$
        \Vert \theta_{q'} - \theta_{q'-1} \Vert_{\dot{B}^\alpha_{\infty,\infty}} + \Vert \theta_{q'} - \theta_{q'-1} \Vert_{L^p} \lesssim 2^{-C(\alpha,p)q'}
        $$
        where the implicit constant depends on $\alpha$ and $p$ but not $q$, or $q'$;
        \item\label{i:5} There is $\delta > 0$ independent of $q$ such that
        $$
        \Vert \theta_q \Vert_{L^1} > (1 + 2^{-q})\delta;
        $$
        \item\label{i:6} For each $q' \leq q$ there exists a unique $j \in \Z$ such that
        $$
        \mathbb{P}_{2^j}(\theta_{q'} - \theta_{q'-1}) = \theta_{q'} - \theta_{q'-1};
        $$
        \item\label{i:7} There are $C_1,C_2 > 0$ independent of $q$ such that
        $$
        \sum_{\substack{n,m \leq q\\ n \not  = m}} \Vert (\theta_n - \theta_{n-1}) (u_m - u_{m-1}) \Vert_{\dot{H}^{-s}} < C_1 - 2^{-q}
        $$
        and
        $$
         \sum_{n \leq q} \Vert (\theta_n - \theta_{n-1}) (u_n - u_{n-1}) \Vert_{\dot{H}^{-s}} < C_2 - 2^{-q + 100}
        $$
    \end{enumerate}
\end{proposition}

\begin{proof}[Proof of Theorem~\ref{thm:main} using Proposition~\ref{prop:ind}]
We start first by verifying the base case. For some $A > 0$ put $\theta_0 = A\cos(2\pi x_1)$, $u_0 = AT(\cos(2\pi x_1))$, and
$$
R_0 = \theta_0 u_0 + \Lambda^{\gamma - 2} \nabla \theta_0.
$$
Items \ref{i:1}, \ref{i:2}, and \ref{i:6} (taking $\theta_{-1} = 0$) are obvious. Taking $A > 0$ small enough ensures items \ref{i:3} and \ref{i:4}. Taking $\delta = A/2$ gives item \ref{i:5}. Choosing $C_1$ and $C_2$ large enough gives item \ref{i:7}.

Now assume all items hold for all $q$. Items \ref{i:1} and \ref{i:4} posit that $\{\theta_q\}$ is a Cauchy sequence in both $\dot{B}^\alpha_{\infty,\infty}(\T^d)$ and $L^p(\T^d)$ for all $\alpha < 0$ and all $p < 2$, and thus has a limit $\theta \in \dot{B}^\alpha_{\infty,\infty}(\T^d) \cap L^p(\T^d)$ for all $\alpha < 0$ and $p < 2$. This gives
$$
\theta \in \bigcap_{0 < \epsilon < 1} \dot{B}^{-\epsilon}_{\infty,\infty}(\T^d) \cap L^{2-\epsilon}(\T^d).
$$
From standard boundedness properties of Calderon-Zygmund operators we conclude $u$ lies in the same space as $\theta$.

From \ref{i:2}, we have that
$$
\int_{\T^d} \nabla \cdot (\theta_q u_q) \psi + \int_{\T^d} \psi \Lambda^\gamma \theta_q = \int_{\T^d} \psi \nabla \cdot R_q.
$$
Applying integration by parts and the self-adjointness of the fractional Laplacian gives
\begin{equation}\label{eq:integral_form}
    -\int_{\T^d} \mathbb{P}_{= 0}(\theta_q u_q) \cdot \nabla \psi + \int_{\T^d} \theta_q \Lambda^\gamma \psi = - \int_{\T^d} R_q \cdot \nabla \psi.
\end{equation}
We want to show that as $q \to \infty$ we obtain a solution in the sense of Definition \ref{def:weak_para_soln}. We start by analyzing the right hand side of ~\eqref{eq:integral_form}. We have
\begin{equation}\label{eq:stress_est}
    \left|\int_{\T^d} R_q \cdot \nabla \psi\right| = \left|\left\langle R_q, \nabla \psi\right\rangle_{\dot{H}^{-s},\dot{H}^s}\right| \lesssim \Vert R_q \Vert_{\dot{H}^{-s}} \Vert \nabla \psi \Vert_{\dot{H}^s} \lesssim 2^{-q} \Vert \nabla \psi \Vert_{\dot{H}^s} \to 0
\end{equation}
as we send $q \to \infty$. Now for the left hand side of Equation ~\eqref{eq:integral_form}, since $\theta_q \to \theta$ in $L^1$ we have
\begin{equation}\label{eq:L1_conv}
    \int_{\T^d} \theta_q \Lambda^\gamma \psi \to \int_{\T^d} \theta \Lambda^\gamma \psi = \langle \theta, \Lambda^\gamma \psi\rangle_{\dot{H}^{-\epsilon},\dot{H}^{\epsilon}}.
\end{equation}
Items \ref{i:6} and \ref{i:7} give that $\mathbb{P}_{\not=0}(\theta u)$ is well defined as a paraproduct in $\dot{H}^{-s}$. Hence
\begin{equation}\label{eq:nonlin_term}
    \begin{split}
        \int_{\T^d} \mathbb{P}_{= 0}(\theta_q u_q) \cdot \nabla \psi &= \int_{\T^d} \nabla \psi \cdot \sum_{\substack{n,m \leq q\\ j,j' \geq 0}} \mathbb{P}_{\not = 0}\left(\mathbb{P}_{2^j}(\theta_n - \theta_{n-1}) \mathbb{P}_{2^{j'}}(u_{m} - u_{m-1})\right)\\
        &= \left\langle \sum_{\substack{n,m \leq q\\ j,j' \geq 0}} \mathbb{P}_{\not = 0}\left(\mathbb{P}_{2^j}(\theta_n - \theta_{n-1}) \mathbb{P}_{2^{j'}}(u_{m} - u_{m-1})\right), \nabla \psi\right\rangle_{\dot{H}^{-s},\dot{H}^s}\\
        &\to \langle \mathbb{P}_{\not=0}(\theta u), \nabla \psi\rangle_{\dot{H}^{-s},\dot{H}^s}
    \end{split}
\end{equation}
Notice the third equality in ~\eqref{eq:nonlin_term} is defined this way using Definition \ref{def:paras}. Hence combining Equations ~\eqref{eq:stress_est}, ~\eqref{eq:L1_conv}, and ~\eqref{eq:nonlin_term} shows $\theta$ and $u$ is a weak solution to ~\eqref{eq:stat_eqn} in the sense of Definition \ref{def:weak_para_soln}; it remains to show that it is nontrivial. And for this, from item \ref{i:5} we see immediately that $\theta \not =0$. If $u = 0$, then we would have
$$
\int_{\T^d} \theta \Lambda^\gamma \psi = 0
$$
for all $\psi$. This forces $\theta = 0$ which is a contradiction. So $u \not =0$, and thus $\theta$ and $u$ are our desired solution.

\end{proof}

\end{document}
